% Source: https://artofproblemsolving.com/community/c3h1577551_series_questions
% Stable identity: aops-topic-1577551-q01
\textbf{(a)} Take $x=2$. The series is $\dfrac{1}{2}+\dfrac{1}{4}+\dfrac{1}{8}+\cdots$, whose partial sums are $1-\dfrac{1}{2^n}$, so its sum is $1$. But $\dfrac{1}{x+1}=\dfrac{1}{3}\ne1$. So the claim is false.

\textbf{(b)} The series is geometric with first term $\dfrac{1}{x}$ and ratio $r=\dfrac{1}{x}$. A geometric series with nonzero first term converges exactly when $|r|<1$, that is when $|x|>1$, and then its sum is
\[
\dfrac{\text{first term}}{1-r}=\dfrac{\dfrac{1}{x}}{1-\dfrac{1}{x}}=\boxed{\dfrac{1}{x-1}}\qquad(|x|>1).
\]
For $|x|\le1$ the terms do not tend to $0$, so the series diverges. (In particular the claimed formula is wrong for every $x$ where the series converges, since $\dfrac{1}{x-1}\ne\dfrac{1}{x+1}$.)

\textbf{(c)} Alternate the signs: the ratio becomes $-\dfrac{1}{x}$, and for $|x|>1$
\[
\dfrac{1}{x}-\dfrac{1}{x^2}+\dfrac{1}{x^3}-\cdots=\dfrac{\dfrac{1}{x}}{1+\dfrac{1}{x}}=\dfrac{1}{x+1}.
\]
So the student's formula holds for the \emph{alternating} series $\displaystyle\sum_{n=1}^\infty\dfrac{(-1)^{n+1}}{x^n}$ with $|x|>1$.
